\paragraph{Edge and contact terms of the vector and connection sector at the Weyl coefficient (conditional; not closure).}
Keep the setting of the preceding paragraph: the weighted \(a_2\) density, the modes with \(\delta\rho=0\) about a constant self-dual or anti-self-dual \(\bar\Theta\), the completed replica, a totally geodesic \(\Sigma\), cross terms at first order in \(1-n\), and the completed gap \(2\kappa_X\) in the coefficient of \(W\), read in the basis \((W,R_\Sigma,R,e,X')\). On the twelve modes \(W\) and \(R_\Sigma\) are independent (rank 2), with \(R_{ii}=W-R_\Sigma\) and \(R_{ijij}=2W-R_\Sigma\).

\emph{Endomorphism contacts.} To first order in the amplitude about a constant \(\bar\Theta\) the endomorphisms \(E_V\) of the vector and \(E_f\) of the Dirac operator equal \(e\,\mathbb I\): the antisymmetric part is a divergence of \(\Theta^{\rho\sigma}\partial_\rho\Theta^{\mu\nu}\), and the symmetric traceless part is quadratic (exact arithmetic, cubic jets of both chiralities, two amplitudes). Since \(e\) vanishes for \(\delta\rho=0\), every cross term of \(a_2\) with a delta of an endomorphism, uniform or confined to the normal block, vanishes at this order for every content: a contact of Kabat type has no coefficient of \(W\) in this presentation. For the Hodge operator of the metric \(G\) (\(E=-\mathrm{Ric}\), \(\Omega=\mathrm{Riem}\)), the same cross-term rule on the same modes gives \(6/5\) and \(-31/15\) for the Maxwell field with its ghosts, equal to its \(12\kappa_C\) in the coefficient of \(W\), of which the endomorphism contact carries \(3\) (the value is \(-9/5\) without it). The matrix-derived vector operator is not the Hodge operator: on these modes \(\mathrm{Riem}^2=2C^2\) and \(\mathrm{Ric}^2=C^2/2\), and \(v-4s\) is \(-C^2/3\) against \(C^2/12\). That mechanism is therefore not available for it.

\emph{Edge modes.} In the classes treated in~\cite{DonnellyWall2016}, products of a two-dimensional base with a bifurcate horizon and a compact fibre, and \(S^1\times H^3\), the edge sector of the Maxwell field is a scalar of wrong-sign statistics on the entangling surface, and its logarithmic density is the two-dimensional trace anomaly of the Laplacian of \(\Sigma\), \(-R_\Sigma/24\pi\); that paper leaves the dependence on the Weyl tensor and on the extrinsic geometry open. In the basis above \(-R_\Sigma\) has coefficient \(0\) of \(W\). A weight \(w\) built from \(\operatorname{Pf}\Theta\) and \(\operatorname{tr}\Theta^2\) in \(-w^{-1}\operatorname{div}(w\nabla)\) adds nothing at first order, because the Hessians of both vanish on the modes. A potential \(E_\Sigma=\alpha R_{ii}+\beta R_{ijij}\) would change the density to \(-(R_\Sigma+6E_\Sigma)\) and would need \(\alpha+2\beta=\kappa_X/3\) for the coefficient \(-2\kappa_X\) of \(W\). This depends on \(n_D\) through the Dirac weight, whereas the vector carries one edge scalar, if it carries the edge sector of a Maxwell field, and the Dirac fields none. No edge term independent of \(n_D\) equals \(2\kappa_X\) for all \(n_D\).

\emph{Holonomy contact.} For the diagonal profile with \(\lvert\Theta^{34}\rvert=r^{-\epsilon c_b}\) the holonomy is \(H=c_a\mathcal J_{12}+\chi c_b\mathcal J_{34}\), with \(\chi=+1\) about a self-dual and \(\chi=-1\) about an anti-self-dual \(\bar\Theta\) (symbolic computation), and the coefficient of \(W\) in \(\operatorname{tr}(\mathcal J_{34}\Omega_{12})\) is \(2\chi\). The gap, in the extended cross-term rule stated above, is therefore the same for both chiralities, \(N_{\mathrm{tot}}(3c_a-2c_b-3)/30+2\kappa_X(2-c_a+c_b)\). With \(s=c_a-c_b\), the contraction of the first-order Weyl tensor with the Weyl tensor of the deltas is \(2sW\), and the Euler combination \(\mathrm{Riem}\cdot\mathrm{Riem}'-4\,\mathrm{Ric}\cdot\mathrm{Ric}'\) is \(2c_aR_\Sigma-2c_bW\), which contains \(W\) for \(c_b\neq0\). If that \(W\) part is dropped as topological, the gap is \((s-1)N_{\mathrm{tot}}/10-2(s-2)\kappa_X\): the metric part requires \(s=1\) and the connection part \(s=2\). The profile with \(\delta\rho=0\), \(c_b=-c_a\), has \(s=2\) and removes the \(\kappa_X\) part, leaving \(N_{\mathrm{tot}}/15\) (\(N_{\mathrm{tot}}/10\) after dropping the Euler part); its induced metric on \(\Sigma\) is singular at first order and the scalar of its Riemann delta, \(2c_a+4c_b\), has the sign opposite to that of the cone.

No edge, endomorphism-contact or holonomy term of this family therefore supplies \(-2\kappa_X W\) for all contents. The required term is proportional to the index-weighted sum that multiplies \(\operatorname{tr}\Omega^2\), the same for every fibre representation; of the candidates examined only a holonomy generator beyond \(\mathcal J_{12}\) attached to \(\Sigma\) has this property, and none of the profiles of the diagonal family provides it without a singular metric on \(\Sigma\).

These statements are for the Poisson-limit operators, one loop, \(\alpha=1\), a totally geodesic \(\Sigma\), cross terms at first order in \(1-n\) on the \(n<1\) branch, the modes and replica of the preceding paragraphs, the cross-term rule extended to \(c_b\neq0\) (not derived), and an edge sector of the abelian type of the cited paper. They do not treat non-diagonal profiles, the non-delta singular first-order parts, the exact heat trace of the cone with holonomy, the symplectic gauge group, or the mode with \(\delta\rho\neq0\), and give no value or sign of \(G_{\mathrm{eff}}\). Problem~B stays open and PROGRAMME-CONDITIONAL.
